% Source: https://artofproblemsolving.com/community/c3h2135932_cutting_a_cake
% Stable identity: aops-topic-2135932-q01
A cake is to be shared among several people. Each person has her own way of valuing pieces of the cake: every piece has a nonnegative value; the values of non-overlapping pieces add up; the whole cake has value $1$; and as a knife is moved continuously across the cake (or a cut is rotated about the centre of a round cake), the value of the piece on one side changes continuously. With $n$ people, a share is \emph{fair} for a person if she values it at least $\dfrac{1}{n}$.
\begin{enumerate}
\item[(a)] Describe a procedure by which two people can divide the cake so that each gets a fair share by his or her own valuation, whatever the other does, and prove that it works.
\item[(b)] Describe such a procedure for $n$ people, and prove that it works.
\item[(c)] For a round cake, the following method has been proposed: each of the $n$ people in turn makes one straight cut from the edge to the centre, and then the people choose pieces in the order in which they cut, the first cutter choosing first. Show that the method guarantees a fair share to each person when $n=2$, but not when $n=3$.
\end{enumerate}
